% Image credits for the photographs and AI illustrations of Book 3 (University, Year 1), Hindi edition.
% Machine-readable license records live in images/book3/CREDITS.md; keep this
% file in sync with frontmatter/image-credits-book3.tex (one \item per item there).
% Names of people, institutions, collections, works and licences are kept
% exactly as in the English credits (Latin script); only the frame and the
% descriptions are translated.
\chapter*{चित्रों का श्रेय}

इस पुस्तक के रेखाचित्र, मानचित्र और चार्ट मौलिक हैं; प्रकाशित आँकड़ों से बनाए गए मानचित्र और
चार्ट अपना स्रोत चित्र-परिचय में या इस सूची में बताते हैं। फ़ोटो उनके अपने-अपने मुक्त लाइसेंसों के
अंतर्गत, उनके लेखकों के प्रति आभार के साथ, इस्तेमाल किए गए हैं; हर फ़ोटो के पूरे स्रोत-लिंक और
लाइसेंस के पते पुस्तक के भंडार में \texttt{images/book3/CREDITS.md} में दिए गए हैं।

\bigskip
फ़ोटो और मौलिक रेखाचित्रों के साथ-साथ, कुछ चित्र AI से बनाए गए चित्रण हैं, जो पुस्तक के
संपादकों के लिखे निर्देशों से OpenAI के चित्र-निर्माण से बनाए गए, और शामिल करने से पहले उनकी
वैज्ञानिक शुद्धता जाँची गई। वे केवल सामान्य दृश्य दिखाते हैं, कभी कोई असली जगह, असली व्यक्ति
या असली नमूना नहीं। हर बनाए गए चित्र का पूरा निर्देश भंडार में
\texttt{images/book3/ai/PROMPTS.md} में दर्ज है।

\bigskip
\noindent फ़ोटो और आँकड़े, अध्याय के अनुसार:
\begin{itemize}\small\setlength{\emergencystretch}{2em}
  \item अध्याय~1 (गोलीय खगोल विज्ञान): Paranal वेधशाला के ऊपर तारों के पथ, R.~Wesson/ESO, CC BY 4.0।
  \item अध्याय~2 (दूरबीनें, संसूचक और प्रकाशमिति): James Webb अंतरिक्ष दूरबीन के प्राथमिक दर्पण के छह खंड, क्रायोजेनिक परीक्षण से पहले, NASA/MSFC/David Higginbotham, सार्वजनिक डोमेन।
  \item अध्याय~3 (तुलनात्मक ग्रह विज्ञान): पार्थिव ग्रह अपने सही अनुपात में, NASA; Magellan रडार से शुक्र, NASA/\allowbreak{}JPL; Viking 1 यान के कक्षीय चित्र से मंगल, NASA/\allowbreak{}USGS; यूरोपा, NASA/\allowbreak{}JPL/\allowbreak{}DLR; आयो, NASA/\allowbreak{}JPL/\allowbreak{}University of Arizona; सभी सार्वजनिक डोमेन।
  \item अध्याय~4 (भूवैज्ञानिक मानचित्र): William Smith, \emph{A Delineation of the Strata of England and Wales, with part of Scotland}, 1815; सार्वजनिक डोमेन (Wikimedia Commons)।
  \item अध्याय~5 (अनुप्रस्थ काट): Siccar Point पर हटन की असंगति, PawełS, CC BY-SA 4.0 (Wikimedia Commons)।
  \item अध्याय~6 (क्षेत्र और प्रयोगशाला विधियाँ I): आड़े ध्रुवकों के बीच देखी गई एक बायोटाइट ग्रेनाइट की पतली काट, Darla Sondrol, CC0 (Wikimedia Commons)।
  \item अध्याय~7 (मुख्य शैल-निर्माता खनिज): ओलिवीन, क्वार्ट्ज़, बायोटाइट, आल्मंडाइन और हॉर्नब्लेंड के नमूने, Robert M.~Lavinsky, CC BY-SA 3.0; लैब्रेडोराइट, Twyla Baker, CC0; कैल्साइट के विदलन समचतुर्षस्रक, Giovanna Canu और Eva Santini, CC BY 4.0 (सभी Wikimedia Commons)।
  \item अध्याय~8 (प्रकाशिकी खनिजविज्ञान): एक गैब्रो की पतली काट, एक ध्रुवक में और आड़े ध्रुवकों के बीच, Chmee2, CC BY-SA 3.0; छात्र के रूप में Auguste Michel-Lévy, सार्वजनिक डोमेन (Wikimedia Commons)। व्यतिकरण रंग CIE 1931 वर्ण-मिलान फलनों और CIE मानक प्रकाश D65 से परिकलित हैं, \copyright{} CIE, CC BY-SA 4.0।
  \item अध्याय~9 (स्तरित पृथ्वी): Andrija Mohorovičić, सार्वजनिक डोमेन; 1932 में Inge Lehmann, Even Neuhaus का छायाचित्र, सार्वजनिक डोमेन (Wikimedia Commons)। PREM: A.~M.~Dziewonski और D.~L.~Anderson (1981), ObsPy में सारणीबद्ध।
  \item अध्याय~10 (पृथ्वी का ऊष्मा-बजट): वैश्विक ऊष्मा प्रवाह J.~H.~Davies और D.~R.~Davies (2010), \emph{Solid Earth}, CC BY 3.0 से; मैंटल का रुद्धोष्म वक्र T.~Katsura (2022), \emph{JGR Solid Earth}, CC BY से।
  \item अध्याय~11 (समस्थिति और गुरुत्व विसंगतियाँ): Vening Meinesz का लोलक उपकरण, Broodt, CC BY-SA 4.0 (Wikimedia Commons)। ऊँचाइयाँ NOAA National Centers for Environmental Information के ETOPO1 से, एक-एक डिग्री पर लिए गए।
  \item अध्याय~12 (गोले पर प्लेट गतिकी): एक स्थायी उपग्रह-स्थिति स्टेशन, Alpsdake, CC BY-SA 4.0 (Wikimedia Commons)। प्लेट सीमाएँ PB2002 से (P.~Bird, 2003; H.~Ahlenius का GeoJSON, ODC-By 1.0); तटरेखाएँ Natural Earth से (सार्वजनिक डोमेन)।
  \item अध्याय~13 (दरारें, कटक और निष्क्रिय किनारे): Þingvellir, आइसलैंड में खुली दरारें, Olga Ernst, CC BY-SA 4.0; Galápagos दरार पर तकिया लावा, NOAA Okeanos Explorer Program, CC BY 2.0 (दोनों Wikimedia Commons)। गहराइयाँ NOAA National Centers for Environmental Information के ETOPO1 से, और समुद्र-तल की आयु M.~Seton और अन्य (2020), EarthByte, CC BY 4.0 से।
  \item अध्याय~14 (सक्रिय किनारे): International Space Station से दिखता Sarychev शिखर का उद्गार, NASA, सार्वजनिक डोमेन (Wikimedia Commons)। भूकंप USGS ANSS Comprehensive Catalog से (सार्वजनिक डोमेन); ज्वालामुखी Smithsonian Institution के Global Volcanism Program से; प्लेट सीमाएँ PB2002 से (P.~Bird, 2003; H.~Ahlenius का GeoJSON, ODC-By 1.0); तटरेखाएँ Natural Earth से (सार्वजनिक डोमेन); समुद्र-तल की आयु M.~Seton और अन्य (2020), EarthByte, CC BY 4.0 से।
  \item अध्याय~15 (हॉटस्पॉट और अंतःप्लेट ज्वालामुखी): Pu\textquoteleft u \textquoteleft O\textquoteleft o, Kilauea पर लावा का फ़व्वारा, G.~E.~Ulrich, U.S. Geological Survey, सार्वजनिक डोमेन; Matheran में पश्चिमी घाट, Nicholas (Nichalp), CC BY-SA 2.5 (दोनों Wikimedia Commons)। हॉटस्पॉट की स्थितियाँ, ट्रैक के दिगंशकोण और दरें W.~J.~Morgan और J.~P.~Morgan (2007) के पूरक की तालिका~1 से; गहराइयाँ NOAA National Centers for Environmental Information के ETOPO1 से; प्लेट सीमाएँ PB2002 से (P.~Bird, 2003; H.~Ahlenius का GeoJSON, ODC-By 1.0); तटरेखाएँ Natural Earth से (सार्वजनिक डोमेन)।
  \item अध्याय~16 (प्रतिबल, विकृति और रियोलॉजी): एक उच्च-दाब त्रिअक्षीय कोष्ठ, जिससे चट्टान की सामर्थ्य मापी जाती है, Darryl Peroni, CSIRO, CC BY 3.0 (Wikimedia Commons)। गीले क्वार्ट्ज़ाइट (Gleason और Tullis, 1995) और शुष्क ओलिवीन (Hirth और Kohlstedt, 2003) के विसर्पण-प्राचल, जैसे M.~A.~Lopez-Sanchez ने \emph{strength\_envelopes} में संकलित किए (MPL-2.0)।
  \item अध्याय~17 (भ्रंश, वलन और विन्यास): Corona Heights, San Francisco का एक स्लिकनसाइड, Dextersinister, CC BY-SA 4.0; Keystone Thrust, Nevada की भ्रंश-ब्रेक्सिया, Jstuby, सार्वजनिक डोमेन; Millook, Cornwall के शेवरॉन वलन, Smalljim, CC BY 3.0; Potes, Cantabria के पास अक्षीय-तलीय विदलन वाली एक अपनति, Mikenorton, CC BY-SA 3.0 (सभी Wikimedia Commons)।
  \item अध्याय~18 (भूकंप का उद्गम): Marin County में 1906 की भ्रंश-दरार से अलग हुई ज़मीन और बाड़ें, \emph{Popular Science Monthly} खंड~69 (1906), सार्वजनिक डोमेन; San Andreas भ्रंश के आर-पार Wallace Creek का विस्थापन, R.~E.~Wallace, U.S. Geological Survey, सार्वजनिक डोमेन (दोनों Wikimedia Commons)। भूकंप USGS ANSS Comprehensive Catalog से (सार्वजनिक डोमेन); 2011 के तोहोकु भूकंप की भ्रंश-क्रियाविधि USGS के W-phase आघूर्ण-टेंसर से (सार्वजनिक डोमेन)।
  \item अध्याय~19 (अपक्षय और मृत्तिकाएँ): Targassonne, Pyrenees में सड़ा ग्रेनाइट और ग्रेनाइट की रेत, Pinpin, CC~BY~2.5; Weipa, Australia में कायोलिनिक बलुआ पत्थर पर बॉक्साइट, Werner Schellmann, CC~BY-SA~2.5; स्कैनिंग इलेक्ट्रॉन सूक्ष्मदर्शी में कायोलिनाइट के क्रिस्टल, U.S. Geological Survey, सार्वजनिक डोमेन (सभी Wikimedia Commons)। आयनिक त्रिज्याएँ R.~D.~Shannon (1976) से।
  \item अध्याय~20 (अपरदन और परिवहन): Terra उपग्रह के MODIS उपकरण से 4~मार्च 2018 को देखी गई Mississippi की बाढ़-प्लूम, NASA Earth Observatory का चित्र, Joshua Stevens द्वारा, सार्वजनिक डोमेन (Wikimedia Commons)।
  \item अध्याय~21 (खंडजात तलछट और उनके परिवेश): Nevada में Aztec Sand\-stone का तिर्यक स्तरण, Chip Carroon, Bureau of Land Management, CC BY-SA~2.0; Landsat~7 से Lena नदी का डेल्टा, U.S. Geological Survey, सार्वजनिक डोमेन; Bouma अनुक्रम के खंडों वाला एक टर्बिडाइट, Becke-Oese, Jo Weber, CC BY-SA~3.0 (सभी Wikimedia Commons)।
  \item अध्याय~22 (कार्बोनेट, सिलिसियस और वाष्पनिक तलछट): अंतरिक्ष स्टेशन से देखा गया Bora Bora, NASA Johnson Space Center, सार्वजनिक डोमेन; \emph{Emiliania huxleyi} की एक कोशिका, Alison R.~Taylor, University of North Carolina Wilmington Microscopy Facility, CC BY~2.5; \emph{Orbulina universa}, H.~Takagi, K.~Kimoto, T.~Fujiki, H.~Saito, C.~Schmidt, M.~Kucera और K.~Moriya, CC BY~4.0; McMurdo Sound के डायटम, Gordon T.~Taylor, Stony Brook University, NOAA Corps Collection, सार्वजनिक डोमेन (सभी Wikimedia Commons)।
  \item अध्याय~23 (डायाजेनेसिस और अवसादी द्रोणियाँ): Potsdam Sand\-stone में क्वार्ट्ज़ की अतिवृद्धियाँ, M.~C.~Rygel, CC BY-SA~4.0 (Wikimedia Commons)। समुद्र-स्तर स्टैक R.~M.~Spratt और L.~E.~Lisiecki (2016), NOAA NCEI World Data Service for Paleo\-climatology से; संपीड़न के प्राचल J.~G.~Sclater और P.~A.~F.~Christie (1980) से, जैसे pyBacktrack (EarthByte) में सारणीबद्ध हैं।
  \item अध्याय~24 (स्तरिक पैमाना): Zumaia में Selandian GSSP का स्वर्ण-स्पाइक, Josugoni, CC BY-SA~4.0 (Wikimedia Commons); Trinidad Lake State Park पर क्रिटेशस--पैलियोजीन सीमा, Nationalparks, CC BY-SA~2.5 (Wikimedia Commons)। International Chronostratigraphic Chart के आधार-आयु और रंग, International Commission on Stratigraphy, CC BY~4.0।
  \item अध्याय~25 (रेडियोमेट्रिक आयु-निर्धारण की विधियाँ): कैथोडोल्यूमिनेसेंस में एक ज़रकॉन-कण, Emmanuel Roquette, CC BY-SA~4.0; एक थर्मल आयनन द्रव्यमान स्पेक्ट्रोमीटर, Radiogenic, CC0 (दोनों Wikimedia Commons)। IntCal20 अंशांकन वक्र, P.~J.~Reimer और अन्य (2020), \emph{Radiocarbon} 62, CC BY~4.0, rintcal पैकेज (M.~Blaauw) के साथ वितरित।
  \item अध्याय~26 (वायुमंडल और महासागर की संरचना): Bering Strait में CTD और रोज़ेट उतारे जा रहे हैं, NASA Goddard Space Flight Center, सार्वजनिक डोमेन (Wikimedia Commons)। World Ocean Atlas 2023 का तापमान और लवणता (NOAA National Centers for Environmental Information)।
  \item अध्याय~27 (विकिरण और ऊष्मा का परिवहन): Deep Space Climate Observatory से सूर्य-प्रकाशित पृथ्वी, NASA Goddard Space Flight Center, सार्वजनिक डोमेन (Wikimedia Commons)। NCEP-NCAR Reanalysis~1 (E.~Kalnay और अन्य, 1996) से विकिरण, सतही फ़्लक्स और वर्षण; NOAA PSL, Boulder, Colorado, USA द्वारा उपलब्ध कराए गए, उनकी वेबसाइट https://psl.noaa.gov से।
  \item अध्याय~28 (पृथ्वी के वायुमंडल का सामान्य परिसंचरण): Gulf of Alaska में एक बहिरुष्णकटिबंधीय चक्रवात, NASA MODIS/LANCE, सार्वजनिक डोमेन; GOES-11 के बादल-आँकड़ों से अंतरोष्ण अभिसरण क्षेत्र, NASA Earth Observatory, सार्वजनिक डोमेन (दोनों Wikimedia Commons)। NCEP-NCAR Reanalysis~1 (E.~Kalnay और अन्य, 1996) से पवनें और समुद्र-तल दाब; NOAA PSL, Boulder, Colorado, USA द्वारा उपलब्ध कराए गए, उनकी वेबसाइट https://psl.noaa.gov से।
  \item अध्याय~29 (महासागर का परिसंचरण): समुद्र-पृष्ठ तापमान में दिखती खाड़ी धारा, NOAA NESDIS, सार्वजनिक डोमेन; तैनात किया जा रहा एक गहरा प्रोफ़ाइलिंग फ़्लोट, अज्ञात NOAA छायाकार, सार्वजनिक डोमेन (दोनों Wikimedia Commons)।
\end{itemize}
